\documentclass[10pt,a4paper]{amsart}
\usepackage[margin=19mm]{geometry}
\usepackage{amsmath,amssymb,mathtools}
\usepackage[scr=rsfs,cal=euler]{mathalfa}
\usepackage{xfrac}
\usepackage[unicode,hidelinks,bookmarks=false]{hyperref}
\newcommand{\ie}{i.e.\ }
\newcommand{\cf}{cf.\ }
\newcommand{\resp}{respectively\ }
\newcommand{\Subsection}{\subsection}
\providecommand{\nobreakdash}{\nobreak\hbox{-}}
\setlength{\emergencystretch}{2em}
\multlinegap0pt
\usepackage{bm}
\usepackage{bbm}
\usepackage{dsfont}
\usepackage{xspace}
\usepackage{cancel}
\usepackage{mathtools}
\usepackage{amsthm}
\makeatletter
\@ifundefined{theorem}{\newtheorem{theorem}{Theorem}}{}
\@ifundefined{corollary}{\newtheorem{corollary}[theorem]{Corollary}}{}
\makeatother
\makeatletter
\expandafter\ifx\csname revdw\endcsname\relax
\newenvironment{revdw}{\color{blue}}{}
\fi
\expandafter\ifx\csname revsg\endcsname\relax
\newenvironment{revsg}{\color{red}}{}
\fi
\expandafter\ifx\csname revvy\endcsname\relax
\newenvironment{revvy}{\color{magenta}}{}
\fi
\makeatother
\title[Borel Combinatorics of Schreier Graphs of $\mathbb{Z}$-actio]{Borel Combinatorics of Schreier Graphs of $\mathbb{Z}$-actions}
\author{Su Gao, Yingying Jiang, Tianhao Wang}
\date{Source: arXiv:2510.27290v1 (CC BY 4.0)}
\begin{document}
\maketitle

\noindent\emph{Source and licence.} Borel Combinatorics of Schreier Graphs of $\mathbb{Z}$-actions. Version arXiv:2510.27290v1 (CC BY 4.0), distributed under the Creative Commons Attribution 4.0 International licence, as recorded in the arXiv licence metadata for this exact version and shown on the abstract page https://arxiv.org/abs/2510.27290v1. Full rights notice: licenses/arxiv-2510.27290-CC-BY-4.0.txt.\par\smallskip

\noindent\textit{Editorial scope.} Counted unit: the Theorem whose source label is \texttt{thm:pair} in the article (source0.tex lines 731--733, source label \texttt{thm:pair}), delivered together with the article's own complete original proof of it (source0.tex lines 735--871, one \texttt{proof} environment of 137 lines). No mathematics is written, repaired or re-derived: statement and proof are the article's own text line for line. Required context retained uncounted: cor:odd (lines 664--665); cor:1ton (lines 669--670); thm:2a1 (lines 683--684). Every definition, display and label of those lines is retained unchanged. Omitted from this package and not redistributed: 1--663, 666--668, 671--682, 685--730, 734--734, 872--1061. Nothing inside a retained excerpt is omitted or altered: every retained body is delivered exactly as the article's lines are written, so no omission receipt of any kind accompanies this package. Citations: the retained excerpts carry no citation command at all, so no bibliography entry of the article is transcribed and no reference is redistributed. Reference adaptation: none was needed and none was made; every reference inside the delivered unit resolves inside it, so no unresolved reference or citation is produced. Printed slips kept literal: no printed word, symbol, hypothesis or number of the counted statement or of its proof was altered. Layout and class: the article's own class is \texttt{amsart} with packages including color, amsmath, amsfonts, amssymb, amsthm, faktor, xy, tikz, caption, subfig, cite, multirow, aliascnt, hyperref; the wrapper is the lane's pinned \texttt{amsart} editorial wrapper at 10pt on A4, which restates the theorem-like environments the retained text needs and adds the article's own macro definitions that the retained text needs (none), with extra wrapper packages (bm, bbm, dsfont, xspace, cancel, mathtools). The delivered numbering is the unit's own. Assets: no retained span contains a figure environment, a graphics-inclusion command, a PDF-page inclusion, a table or a TikZ picture; no image file is copied into this package, the package declares no asset dependency and no image file is redistributed with it. Licence: version arXiv:2510.27290v1 of this article is distributed under the Creative Commons Attribution 4.0 International licence, as recorded in the arXiv licence metadata for this exact version and shown on the abstract page https://arxiv.org/abs/2510.27290v1. A full rights notice is delivered at licenses/arxiv-2510.27290-CC-BY-4.0.txt. Characters: every retained excerpt is pure ASCII, so the delivered engine, XeLaTeX with its default Latin Modern text font, reports no missing character.
\begin{corollary}\label{cor:odd} If $(a_1, \dots, a_n)\in \Sigma$ and all of $a_1, \dots, a_n$ are odd numbers, then $\chi(a_1,\dots, a_n)=3$.
\end{corollary}

\begin{corollary}\label{cor:1ton} For any $n\geq 1$, $\chi(1, 2,\dots, n)=n+2$.
\end{corollary}

\begin{theorem}\label{thm:2a1} If $(a_1, \dots, a_n)\in \Sigma$ and $a_n<2a_1$, then $\chi(a_1,\dots, a_n)=3$.
\end{theorem}

\begin{theorem}\label{thm:pair} If $(a_1, a_2)\in \Sigma$, then
$$ \chi(a_1, a_2)=\left\{\begin{array}{ll}4, & \mbox{if $(a_1, a_2)=(1,2)$,} \\3, & \mbox{otherwise.}\end{array}\right. $$
\end{theorem}

\begin{proof} Let $S=\{\pm a_1, \pm a_2\}$. The fact $\chi(1,2)=4$ is by Corollary~\ref{cor:1ton}. We show that whenever $(a_1, a_2)\neq (1,2)$, we have $\chi(a_1, a_2)=3$. By Corollary~\ref{cor:odd}, we have $\chi(1, a_2)=3$ if $a_2$ is odd. Next we deal with the case $a_1=1$ and $a_2$ is even.

Case 1: $a_1=1$ and $a_2$ is even. 

Let $\ell=a_2+1$ and $t=a_2/2-1$. Let $c_0=(01)^t021$. Then $c_0$ has length $\ell$. Let $p=3\ell$ and $q=3\ell-2$. Then $p, q>\ell$ and $\gcd(p,q)=1$. Let $c_1=c_0^3$. By straightforward observation, $c_1$ is an $S$-coloration. To make this observation easier, we list the three copies of $c_0$ in three rows, where in each column, the adjacent entries correspond to positions which differ by $a_2=\ell-1$.
$$ \begin{array}{rccccccccccccl}
0 & 1 & 0 & 1 & \cdots & 0 & 1 & 0 & 1 & 0 & 2 & 1 &  & \\
   & 0 & 1 & 0 & \cdots &    & 0 & 1 & 0 & 1 & 0 & 2 & 1 & \\
   &    & 0 & 1 & \cdots &    &    & 0 & 1 & 0 & 1 & 0 & 2 & 1
   \end{array}
   $$
Let $c_2$ be obtained from $c_1$ by removing the subword $01$ in its $\ell+1$ and $\ell+2$ coordinates. Then $c_2$ is still an $S$-coloration by a straightforward observation of the following diagram.
$$ \begin{array}{rccccccccccccl}
0 & 1 & 0 & 1 & \cdots & 0 & 1 & 0 & 1 & 0 & 2 & 1 &  & \\
   & 0 & 1 & 0 & \cdots &    & 0 & 1 & 0 & 2 & 1 &  &  & \\
0   &  1  & 0 & 1 & \cdots &    &    & 0 & 1 & 0 & 2 & 1 &  &
   \end{array}
   $$

For the rest of the proof, we assume $a_1>1$. By Theorem~\ref{thm:2a1}, we may assume $a_2\geq 2a_1$. Since $(a_1, a_2)\in\Sigma$, $\gcd(a_1, a_2)=1$. Thus we can finish the proof by considering the following two cases. 

Case 2: For an odd number $m>1$, $ma_1<a_2<(m+1)a_1$.

In this part of the proof we redefine the values of $\ell, t, p, q, c_0, c_1$ and $c_2$. Let $b=a_2-ma_1$. Then $1\leq b<a_1$. Let $\ell=a_1+a_2=(m+1)a_1+b$ and $t=(m-1)/2\geq 1$. As before, we first define an $S$-coloration $c_0$ of length $\ell$ as follows. Let $u=0^{a_1}$, $v=1^{a_1-1}2$, $w=2^{a_1}$, and $s=1^b$. Then let
$$ c_0=(uv)^t usw. $$
It is easily seen that $u,v,w$ are of length $a_1$ and $c_0$ is then of length $2a_1(t+1)+b=\ell$. Also note that 
\begin{itemize}
\item $u$ and $v$ differ in every position,
\item $u$ and $w$ differ in every position, and
\item $sw$ and $us$ both have length $a_1+b$ and they differ in every position. 
\end{itemize}
These properties guarantee that $c_0$ is an $S$-coloration. Let $p=3\ell$ and $q=p-1$. Then $p, q>\ell$ and $\gcd(p, q)=1$. Let $c_1=c_0^3$. Then we may again list $c_1$ in three rows where the adjacent entries in a column correspond to positions which differ by $a_2=ma_1+b$, and observe that $c_1$ is an $S$-coloration. In the following diagram, we use square brackets to group $sw$ together in the first row and $us$ together in the second row, and use parentheses to group $sw$ together in the second row and $us$ together in the third row to show their correspondence in position.
$$  \begin{array}{rccccccccccccl}
u & v & u & v & \cdots & u & v & u & v & u & [s & w] &  & \\
   & u & v & u & \cdots &    & u & v & u & v & [u & (s] & w) & \\
   &    & u & v & \cdots &    &    & u & v & u & v & (u & s) & w
   \end{array}
   $$

Before we define $c_2$ we consider the word $c_0^*$, which is obtained from $c_0$ by removing its first entry (the $0$ in the $0$ coordinate). Then we may write $c_0^*$ as
$$ c_0^*=(u^*v^*)^t u^*s^*w^-, $$
where $u^*=0^{a_1-1}1$, $v^*=1^{a_1-2}20$, $s^*=1^{b-1}2$ and $w^-=2^{a_1-1}$. Now define $\tilde{v}=1^{a_1-2}2^2=1^{a_1-2}22$, and let
$$ \tilde{c}_0=(u^*\tilde{v})^tu^*sw^-. $$
Then note that 
\begin{itemize}
\item $u^*$ and $\tilde{v}$ are both of length $a_1$ and they differ in every position, 
\item $w^-$ is of length $a_1-1$ and $u^*$ and $w^-$ differ in every one of the first $a_1-1$ positions, 
\item $sw$ and $u^*s$ are both of length $a_1+b$ and they differ in every position, and
\item $sw^-$ has length $a_1+b-1$ and $sw^-$ and $u^*s$ differ in every one of the first $a_1+b-1$ positions.
\end{itemize}
It follows that $\tilde{c_0}$ is an $S$-coloration. Moreover, we now have that 
\begin{itemize}
\item $v$ and $u^*$ both have length $a_1$ and they differ in every position, and
\item $u$ and $\tilde{v}$ both have length $a_1$ and they differ in every position.  
\end{itemize}
It follows that $c_0\tilde{c_0}$ is an $S$-coloration by the following diagram.
$$  \begin{array}{rccccccccccccl}
u & v & u & v & \cdots & u & v & u & v & u & [s & w] &  & \\
   & u^* & \tilde{v} & u^* & \cdots &    & u^* & \tilde{v} & u^* & \tilde{v} & [u^* & s] & w^- & 
      \end{array}
   $$
Now we may rewrite $\tilde{c}_0$ as
$$ \tilde{c}_0=u^{-}v(\tilde{u}v)^{t-1}\tilde{u}s\tilde{w}, $$
where $u^{-}=0^{a_1-1}$, $\tilde{u}=20^{a_1-1}$ and $\tilde{w}=12^{a_1-1}$. Note that
\begin{itemize}
\item $\tilde{u}$ and $v$ both have length $a_1$ and they differ in every position, and
\item $s\tilde{w}$ and $us$ both have length $a_1+b$ and they differ in every position.
\end{itemize}
It follows that $\tilde{c}_0c_0$ is also an $S$-coloration by the following diagram.
$$  \begin{array}{rccccccccccccl}
    u^{-} & v& \tilde{u} & \cdots &    & \tilde{u} &v & \tilde{u} & v & \tilde{u} & [s & \tilde{w}] & \\  
& u & v & \cdots &  & & u & v & u & v & [u & s] & w &       \end{array}
   $$
Finally, define $c_2=c_0\tilde{c}_0c_0$. Then $c_2$ is an $S$-coloration of length $q$. This finishes the proof for Case 2.


   
Case 3: For an even number $m>1$, $ma_1<a_2<(m+1)a_1$. 

In this part of the proof we again redefine the values of $\ell, t, p, q, c_0, c_1$ and $c_2$. But we will keep the definitions of $b, u, v, s, u^*, v^*, s^*, \tilde{v}, \tilde{u}$ and $u^{-}$. For the convenience of the reader, we summarize their definitions in the following.
$$\begin{array}{llll}
u=0^{a_1} & u^*=0^{a_1-1}1 & \tilde{u}=20^{a_1-1} & u^-=0^{a_1-1}\\
v=1^{a_1-1}2 & v^*=1^{a_1-2}20 & \tilde{v}=1^{a_1-2}22 & \\
%w=2^{a_1} & & \tilde{w}=12^{a_1-1} & w^-=2^{a_1-1} \\
s=1^b & s^*=1^{b-1}2 & & 
\end{array}
$$
We let $b=a_2-ma_1$ and $d=a_1-b$. Let $t=m/2-1$ and $\ell=a_1+a_2$. Note that
$$\ell=2a_1t+a_1+b+d+b+a_1.$$
Let $p=3\ell$ and $q=3\ell-1$. Then $p, q>\ell$ and $\gcd(p,q)=1$. 

We add the following definitions.
$$ \begin{array}{llll}
x=2^d & x^*=2^{d-1}0 & & \\
y=0^b & y^*=0^{b-1}1 &  & \\
z=1^{a_1} & & \tilde{z}=01^{a_1-1} & z^-=1^{a_1-1}
\end{array}
$$ 
Then let
$$ c_0=(uv)^tusxyz $$
and $c_1=c_0^3$. Then $c_0$ has length $\ell$ and $c_1$ has length $3\ell$. Now note that
\begin{itemize}
\item $u$ and $z$ both have length $a_1$ and they differ in every position,
\item $s$ and $y$ both have length $b$  and they differ in every position,
\item $sx$ and $u$ both have length $a_1$ and they differ in every position, and
\item $xy$ and $z$ both have length $a_1$ and they differ in every position.
\end{itemize}
It follows that both $c_0$ and $c_1$ are $S$-colorations by the following diagram.
$$  \begin{array}{rcccccccccccccl}
    u & v & u & v & \cdots & u & v & u & v & u & sx &  y & z & \\
    & u & v & u & \cdots & & u& v& u & v& u& s & xy & z     \end{array}
   $$
Now let $c_0^*$ be obtained from $c_0$ by removing its first entry. Then we can rewrite $c_0^*$ as
$$ c_0^*=(u^*v^*)^tu^*s^*x^*y^*z^{-}.$$
Now let
$$ \tilde{c}_0=(u^*\tilde{v})^tu^*s^*x^*yz^{-}. $$
Note that
\begin{itemize}
\item $sx$ and $u^*$ still differ in every position, 
\item $s^*$ and $y$ still differ in every position, and
\item $x^*y$ and $z$ still differ in every position.
\end{itemize}
It follows that $c_0\tilde{c}_0$ is an $S$-coloration by the following diagram.
$$  \begin{array}{rcccccccccccccl}
    u & v & u & v & \cdots & u & v & u & v & u & sx &  y & z & \\
    & u^* & \tilde{v} & u^* & \cdots & & u^*& \tilde{v}& u^* & \tilde{v}& u^*& s^* & x^*y & z^{-}    \end{array}
   $$
$\tilde{c}_0$ can be rewritten as 
   $$ \tilde{c}_0=u^{-}v(\tilde{u}v)^{t-1}\tilde{u}sxy\tilde{z}. $$
   Note that $xy$ and $\tilde{z}$ still differ in every position. It follows that $\tilde{c}_0c_0$ is an $S$-coloration by the following diagram.
   $$  \begin{array}{rcccccccccccccl}
   u^{-} & v & \tilde{u} & v& \cdots & v & \tilde{u} & v &\tilde{u} & sx & y & \tilde{z} \\
   & u & v & u & \cdots &  & v & u & v & u & s &  xy & z & \\
    \end{array}
   $$
Now in all cases we have shown that $\chi(a_1, a_2)=3$ if $(a_1, a_2)\in \Sigma$ and $(a_1, a_2)\neq (1,2)$. The theorem is proved.
\end{proof}

\end{document}
